\section{Non-TRM addendum}
This addendum strips TRM out of the non-LLM story and keeps only Compress ARC and VARC. The useful efficiency proxies here are final performance, oracle candidate coverage, selection efficiency, and the pass@k progression. Because Compress ARC has one run and VARC has four, this is descriptive rather than a formal inferential comparison.

\subsection{Hypotheses}
\begin{itemize}
\item Final performance is limited partly by selection, not only candidate generation.
\item ARC-1 should be more selection-efficient than ARC-2.
\item VARC gains should diminish after pass@2.
\end{itemize}

\subsection{Not pursued}
\begin{itemize}
\item We did not pool TRM into this addendum, because step count is not on the same scale as candidate-count or iteration proxies.
\item We did not treat Compress ARC as inferential evidence, because it contributes only a single run.
\item We did not force one unified non-LLM efficiency score across all subtypes, because the telemetry is structurally different across Compress ARC, VARC, and TRM.
\end{itemize}

\subsection{Summary table}
\begin{table}[htbp]
\centering
\caption{Non-TRM efficiency summary.}
\label{tab:nontrm_summary}
\begin{tabular}{lrrrrrr}
\toprule
run & performance & oracle & gap & selection & pass1 & pass4 \\
\midrule
Compress ARC & 0.203 & 0.343 & 0.140 & 0.591 & -- & -- \\
ARC-1\_Unet & 0.367 & 0.623 & 0.255 & 0.590 & 0.300 & 0.367 \\
ARC-1\_ViT & 0.450 & 0.703 & 0.253 & 0.641 & 0.300 & 0.450 \\
ARC-2\_Unet & 0.033 & 0.092 & 0.058 & 0.364 & 0.017 & 0.033 \\
ARC-2\_ViT & 0.042 & 0.125 & 0.083 & 0.333 & 0.025 & 0.042 \\
\bottomrule
\end{tabular}
\end{table}

\subsection{ARC-1 versus ARC-2 within VARC}
Because VARC has two ARC-1 and two ARC-2 runs, we can do an exact two-group permutation check on the split effect. This is still descriptive in spirit, but it gives a clean sanity check for whether ARC-1 is systematically easier to exploit.

\begin{table}[htbp]
\centering
\caption{Exact split comparison for the VARC runs.}
\label{tab:varc_split}
\begin{tabular}{lrrrr}
\toprule
metric & ARC-1 mean & ARC-2 mean & difference & exact p \\
\midrule
performance\_rate & 0.409 & 0.037 & 0.371 & 0.3333 \\
oracle\_rate & 0.663 & 0.108 & 0.554 & 0.3333 \\
selection\_efficiency & 0.615 & 0.348 & 0.267 & 0.3333 \\
pass\_gain\_1\_to\_4 & 0.109 & 0.017 & 0.092 & 0.3333 \\
\bottomrule
\end{tabular}
\end{table}

\subsection{Takeaway}
The main pattern is that oracle coverage is consistently larger than final performance, which points to selection as an important bottleneck. Compress ARC and ARC-1 Unet sit around 0.59 selection efficiency, ARC-1 ViT is a bit better at about 0.64, and ARC-2 drops to roughly 0.33--0.36. That makes the ARC-2 VARC runs look less like a search-depth problem and more like a harder candidate-ranking problem.

\begin{figure}[htbp]
\centering
\includegraphics[width=0.96\linewidth]{figures/nontrm_ceiling_and_selection.png}
\caption{Non-TRM candidate-ceiling view. Left: final performance versus oracle candidate coverage. Right: selection efficiency, defined as final performance divided by oracle rate.}
\label{fig:nontrm_ceiling}
\end{figure}

\begin{figure}[htbp]
\centering
\includegraphics[width=0.86\linewidth]{figures/nontrm_varc_pass_curves.png}
\caption{VARC pass@k curves. Gains are front-loaded and flatten quickly, especially on ARC-2.}
\label{fig:nontrm_passk}
\end{figure}
